% Options for packages loaded elsewhere
% Options for packages loaded elsewhere
\PassOptionsToPackage{unicode}{hyperref}
\PassOptionsToPackage{hyphens}{url}
\PassOptionsToPackage{dvipsnames,svgnames,x11names}{xcolor}
%
\documentclass[
  english,
  letterpaper,
  DIV=11,
  numbers=noendperiod]{scrartcl}
\usepackage{xcolor}
\usepackage{amsmath,amssymb}
\setcounter{secnumdepth}{-\maxdimen} % remove section numbering
\usepackage{iftex}
\ifPDFTeX
  \usepackage[T1]{fontenc}
  \usepackage[utf8]{inputenc}
  \usepackage{textcomp} % provide euro and other symbols
\else % if luatex or xetex
  \usepackage{unicode-math} % this also loads fontspec
  \defaultfontfeatures{Scale=MatchLowercase}
  \defaultfontfeatures[\rmfamily]{Ligatures=TeX,Scale=1}
\fi
\usepackage{lmodern}
\ifPDFTeX\else
  % xetex/luatex font selection
\fi
% Use upquote if available, for straight quotes in verbatim environments
\IfFileExists{upquote.sty}{\usepackage{upquote}}{}
\IfFileExists{microtype.sty}{% use microtype if available
  \usepackage[]{microtype}
  \UseMicrotypeSet[protrusion]{basicmath} % disable protrusion for tt fonts
}{}
\makeatletter
\@ifundefined{KOMAClassName}{% if non-KOMA class
  \IfFileExists{parskip.sty}{%
    \usepackage{parskip}
  }{% else
    \setlength{\parindent}{0pt}
    \setlength{\parskip}{6pt plus 2pt minus 1pt}}
}{% if KOMA class
  \KOMAoptions{parskip=half}}
\makeatother
% Make \paragraph and \subparagraph free-standing
\makeatletter
\ifx\paragraph\undefined\else
  \let\oldparagraph\paragraph
  \renewcommand{\paragraph}{
    \@ifstar
      \xxxParagraphStar
      \xxxParagraphNoStar
  }
  \newcommand{\xxxParagraphStar}[1]{\oldparagraph*{#1}\mbox{}}
  \newcommand{\xxxParagraphNoStar}[1]{\oldparagraph{#1}\mbox{}}
\fi
\ifx\subparagraph\undefined\else
  \let\oldsubparagraph\subparagraph
  \renewcommand{\subparagraph}{
    \@ifstar
      \xxxSubParagraphStar
      \xxxSubParagraphNoStar
  }
  \newcommand{\xxxSubParagraphStar}[1]{\oldsubparagraph*{#1}\mbox{}}
  \newcommand{\xxxSubParagraphNoStar}[1]{\oldsubparagraph{#1}\mbox{}}
\fi
\makeatother


\usepackage{longtable,booktabs,array}
\newcounter{none} % for unnumbered tables
\usepackage{calc} % for calculating minipage widths
% Correct order of tables after \paragraph or \subparagraph
\usepackage{etoolbox}
\makeatletter
\patchcmd\longtable{\par}{\if@noskipsec\mbox{}\fi\par}{}{}
\makeatother
% Allow footnotes in longtable head/foot
\IfFileExists{footnotehyper.sty}{\usepackage{footnotehyper}}{\usepackage{footnote}}
\makesavenoteenv{longtable}
\usepackage{graphicx}
\makeatletter
\newsavebox\pandoc@box
\newcommand*\pandocbounded[1]{% scales image to fit in text height/width
  \sbox\pandoc@box{#1}%
  \Gscale@div\@tempa{\textheight}{\dimexpr\ht\pandoc@box+\dp\pandoc@box\relax}%
  \Gscale@div\@tempb{\linewidth}{\wd\pandoc@box}%
  \ifdim\@tempb\p@<\@tempa\p@\let\@tempa\@tempb\fi% select the smaller of both
  \ifdim\@tempa\p@<\p@\scalebox{\@tempa}{\usebox\pandoc@box}%
  \else\usebox{\pandoc@box}%
  \fi%
}
% Set default figure placement to htbp
\def\fps@figure{htbp}
\makeatother



\ifLuaTeX
\usepackage[bidi=basic,shorthands=off]{babel}
\else
\usepackage[bidi=default,shorthands=off]{babel}
\fi
\ifLuaTeX
  \usepackage{selnolig} % disable illegal ligatures
\fi


\setlength{\emergencystretch}{3em} % prevent overfull lines

\providecommand{\tightlist}{%
  \setlength{\itemsep}{0pt}\setlength{\parskip}{0pt}}



 


\KOMAoption{captions}{tableheading}
\makeatletter
\@ifpackageloaded{caption}{}{\usepackage{caption}}
\AtBeginDocument{%
\ifdefined\contentsname
  \renewcommand*\contentsname{Table of contents}
\else
  \newcommand\contentsname{Table of contents}
\fi
\ifdefined\listfigurename
  \renewcommand*\listfigurename{List of Figures}
\else
  \newcommand\listfigurename{List of Figures}
\fi
\ifdefined\listtablename
  \renewcommand*\listtablename{List of Tables}
\else
  \newcommand\listtablename{List of Tables}
\fi
\ifdefined\figurename
  \renewcommand*\figurename{Figure}
\else
  \newcommand\figurename{Figure}
\fi
\ifdefined\tablename
  \renewcommand*\tablename{Table}
\else
  \newcommand\tablename{Table}
\fi
}
\@ifpackageloaded{float}{}{\usepackage{float}}
\floatstyle{ruled}
\@ifundefined{c@chapter}{\newfloat{codelisting}{h}{lop}}{\newfloat{codelisting}{h}{lop}[chapter]}
\floatname{codelisting}{Listing}
\newcommand*\listoflistings{\listof{codelisting}{List of Listings}}
\makeatother
\makeatletter
\makeatother
\makeatletter
\@ifpackageloaded{caption}{}{\usepackage{caption}}
\@ifpackageloaded{subcaption}{}{\usepackage{subcaption}}
\makeatother
\usepackage{bookmark}
\IfFileExists{xurl.sty}{\usepackage{xurl}}{} % add URL line breaks if available
\urlstyle{same}
\hypersetup{
  pdftitle={Furrballs: A NUMA-Aware{,} High-Performance Data Platform for Low-Latency Server Workloads},
  pdfauthor={NuAtlas},
  pdflang={en},
  colorlinks=true,
  linkcolor={blue},
  filecolor={Maroon},
  citecolor={Blue},
  urlcolor={Blue},
  pdfcreator={LaTeX via pandoc}}


\title{Furrballs: A NUMA-Aware, High-Performance Data Platform for
Low-Latency Server Workloads}
\usepackage{etoolbox}
\makeatletter
\providecommand{\subtitle}[1]{% add subtitle to \maketitle
  \apptocmd{\@title}{\par {\large #1 \par}}{}{}
}
\makeatother
\subtitle{White Paper Draft}
\author{NuAtlas}
\date{2026-04-09}
\begin{document}
\maketitle


\subsection{Status}\label{status}

Draft v0.2 (living document)

\subsection{Abstract}\label{abstract}

Furrballs is a high-performance data platform designed for
latency-sensitive workloads where memory hierarchy behavior, CPU
locality, and throughput consistency are critical. The project combines
a modular cache architecture, compression-aware storage integration, and
an evolving NUMA-aware memory strategy intended to reduce cross-node
penalties and stabilize tail latency under load. Development is
intentionally phased: first a game-engine-focused core library, then a
more feature-complete general library, then a server/client system with
an explicit protocol.

This white paper defines the problem space, design goals, architecture
boundaries, and measurement plan for Furrballs. It also documents
project rescoping history, current implementation status, planned
capabilities, and a reproducible benchmarking approach focused on
p95/p99 latency, throughput, memory efficiency, and NUMA locality
effects.

\subsection{Project Narrative and
Rescoping}\label{project-narrative-and-rescoping}

Furrballs evolved through several scope changes that matter for both
technical context and decision traceability:

\begin{enumerate}
\def\labelenumi{\arabic{enumi}.}
\tightlist
\item
  Initial motivation: investigate NUMA-aware caching after observing
  that popular systems such as Redis are not strongly NUMA-aware for all
  deployment scenarios.
\item
  Practical inspiration: build a cache-first library influenced by
  Ristretto's simplicity and performance-oriented design goals.
\item
  Constraint phase: no direct access to NUMA hardware, so the idea
  expanded into simulating distribution pressure by spreading cache
  behavior across multiple machines.
\item
  Downscope phase: temporary shift toward game-engine use cases (the
  author's strongest domain) to ensure continuous progress and ship a
  usable core.
\item
  Current rescope: return to the larger high-performance objective with
  NUMA-awareness as the primary long-term focus.
\end{enumerate}

This narrative is intentional, not accidental. Each scope change reduced
risk, increased implementation confidence, or unlocked a clearer path to
validation.

\subsection{1. Problem Statement}\label{problem-statement}

Modern low-latency services often degrade under memory pressure,
irregular access patterns, and multi-socket contention. Typical
bottlenecks include:

\begin{itemize}
\tightlist
\item
  Tail-latency spikes from cache misses and allocator fragmentation.
\item
  Performance loss from non-local NUMA memory access.
\item
  Operational complexity in combining cache, compression, and
  persistence layers.
\item
  Difficulty validating behavior consistently across different hardware
  and deployment models.
\end{itemize}

Furrballs targets these bottlenecks with an architecture that
prioritizes locality, predictable behavior, and modular policy design.

\subsection{2. Goals and Non-Goals}\label{goals-and-non-goals}

\subsubsection{Goals}\label{goals}

\begin{itemize}
\tightlist
\item
  Low and stable latency under mixed read/write load.
\item
  NUMA-aware allocation and access behavior.
\item
  Modular cache policy layer with pluggable eviction strategies.
\item
  Progressive delivery through phased architecture evolution.
\item
  Reproducible performance evaluation and transparent tradeoff
  reporting.
\end{itemize}

\subsubsection{Non-Goals (current phase)}\label{non-goals-current-phase}

\begin{itemize}
\tightlist
\item
  General-purpose database replacement.
\item
  Broad feature parity with large data platforms.
\item
  Platform-specific tuning for all operating systems at once.
\end{itemize}

\subsection{3. Design Principles}\label{design-principles}

\begin{itemize}
\tightlist
\item
  Locality first: optimize for CPU and memory locality before adding
  features.
\item
  Explicit tradeoffs: prefer clear, measurable behavior over opaque
  heuristics.
\item
  Incremental validation: every subsystem change should map to benchmark
  deltas.
\item
  Modularity: cache policies, memory strategies, and transport should
  evolve independently.
\end{itemize}

\subsection{4. High-Level Architecture}\label{high-level-architecture}

\subsubsection{Layer boundaries}\label{layer-boundaries}

\begin{itemize}
\tightlist
\item
  Furrballs Core (library): owns memory management, cache policy
  mechanics, compression integration, and storage interactions.
\item
  Furrballs Server (wrapper layer): hosts the active cache runtime,
  wraps Core functionality, and exposes network-facing operations.
\item
  Furrballs Client (interface layer): communicates with the Server
  through protocol-defined requests and does not host cache state.
\end{itemize}

Only the Server hosts and mutates cache state. The Client is
intentionally thin and transport-oriented.

\subsubsection{Core components}\label{core-components}

\begin{itemize}
\tightlist
\item
  Memory layer: page-oriented allocation and management, evolving toward
  NUMA-aware placement.
\item
  Cache layer: policy-driven in-memory caching with extensible eviction
  behavior.
\item
  Storage layer: persistent backing integration with compression
  support.
\item
  Runtime layer: server process, client protocol/API, and operational
  instrumentation.
\end{itemize}

\subsubsection{Data path (conceptual)}\label{data-path-conceptual}

\begin{enumerate}
\def\labelenumi{\arabic{enumi}.}
\tightlist
\item
  Client request arrives at the Server through the protocol boundary.
\item
  Server routes the request through Core cache and locality-aware lookup
  path.
\item
  Miss path loads or reconstructs data from backing store via Core.
\item
  Returned data is inserted/promoted by Server-hosted cache policy and
  memory placement rules.
\item
  Telemetry captures latency, locality, and resource counters.
\end{enumerate}

\subsection{5. NUMA-Aware Direction}\label{numa-aware-direction}

Planned NUMA-awareness work includes:

\begin{itemize}
\tightlist
\item
  Node-aware allocation policies to reduce remote memory traffic.
\item
  Affinity-aware worker scheduling where practical.
\item
  Topology-aware cache partitioning or sharding exploration.
\item
  Local/remote access ratio instrumentation.
\end{itemize}

Key hypothesis: locality-aware placement and scheduling can improve p99
latency and reduce jitter at equivalent throughput.

\subsection{6. Server/Client Expansion
Plan}\label{serverclient-expansion-plan}

\subsubsection{Server responsibilities}\label{server-responsibilities}

\begin{itemize}
\tightlist
\item
  Request handling, cache management, and persistence coordination
  through Core.
\item
  Concurrency control and worker model.
\item
  Metrics and health endpoints.
\end{itemize}

\subsubsection{Client responsibilities}\label{client-responsibilities}

\begin{itemize}
\tightlist
\item
  Connection management and request batching options.
\item
  Retry/backoff strategy (policy-defined).
\item
  Latency and error instrumentation hooks.
\item
  No in-process cache ownership.
\end{itemize}

\subsubsection{Protocol concerns}\label{protocol-concerns}

\begin{itemize}
\tightlist
\item
  Message framing and versioning.
\item
  Backpressure behavior.
\item
  Fault handling and timeout semantics.
\end{itemize}

\subsection{7. Evaluation Methodology}\label{evaluation-methodology}

\subsubsection{7.1 Metrics}\label{metrics}

Primary metrics: - Read latency: p50, p95, p99. - Write latency: p50,
p95, p99. - Throughput: operations per second. - Cache hit ratio. -
Compression ratio and CPU overhead. - Memory overhead and fragmentation
behavior. - NUMA locality indicators (local vs remote memory access
where measurable).

\subsubsection{7.2 Workload Profiles}\label{workload-profiles}

\begin{itemize}
\tightlist
\item
  Read-heavy steady-state.
\item
  Write-heavy burst.
\item
  Mixed workload with skewed key popularity.
\item
  Large-object stress path.
\item
  Multi-client concurrency scaling.
\end{itemize}

\subsubsection{7.3 Benchmark Matrix}\label{benchmark-matrix}

\begin{itemize}
\tightlist
\item
  Single-socket vs multi-socket host.
\item
  NUMA policy variations.
\item
  Cache policy variations.
\item
  Compression on/off.
\item
  Thread count scaling sweep.
\end{itemize}

\subsubsection{7.4 Reproducibility
Requirements}\label{reproducibility-requirements}

\begin{itemize}
\tightlist
\item
  Fixed build configuration per benchmark run.
\item
  Logged hardware and kernel/runtime metadata.
\item
  Versioned benchmark scripts and datasets.
\item
  Raw results retained with timestamp and commit identifier.
\end{itemize}

\subsection{8. Claims and Evidence
Policy}\label{claims-and-evidence-policy}

Any performance claim in this white paper should include: - The exact
workload profile. - The hardware/software test environment. - Comparison
baseline and configuration parity. - Statistical summary from repeated
runs.

\subsection{9. Current Status Snapshot}\label{current-status-snapshot}

Current implementation includes foundational cache/storage abstractions
and early integration points. Development is currently centered on
strengthening Furrballs Core for game-engine-relevant workloads as Phase
1. Several advanced components are still under active development,
especially around full NUMA strategy, server/client protocol maturity,
and end-to-end benchmark coverage.

This white paper is intentionally structured as a living technical
document and should be updated alongside implementation milestones.

\subsection{10. Risks and Open
Questions}\label{risks-and-open-questions}

\begin{itemize}
\tightlist
\item
  Which NUMA strategies provide best p99 gains for your target workload
  mix?
\item
  What is the throughput vs tail-latency tradeoff at high concurrency?
\item
  How should cache policy adapt under rapidly changing access
  distributions?
\item
  What observability granularity is required without adding excessive
  overhead?
\end{itemize}

\subsection{11. Milestone-Driven
Roadmap}\label{milestone-driven-roadmap}

\subsubsection{Phase 1: Furrballs Core for game
engines}\label{phase-1-furrballs-core-for-game-engines}

\begin{itemize}
\tightlist
\item
  Deliver a stable, documented cache/compression/storage core library.
\item
  Validate low-latency behavior on representative game-engine-style
  access patterns.
\item
  Build correctness and regression coverage around core APIs.
\end{itemize}

\subsubsection{Phase 2: Feature-complete library
evolution}\label{phase-2-feature-complete-library-evolution}

\begin{itemize}
\tightlist
\item
  Expand policy surface, observability hooks, and integration
  ergonomics.
\item
  Improve memory strategy and prepare NUMA-related interfaces and
  metrics.
\item
  Broaden benchmark coverage and improve reproducibility on VPS and
  local systems.
\end{itemize}

\subsubsection{Phase 3: Server + Client +
Protocol}\label{phase-3-server-client-protocol}

\begin{itemize}
\tightlist
\item
  Introduce Furrballs Server that wraps Core and hosts cache state.
\item
  Introduce Furrballs Client that interfaces with Server only.
\item
  Define protocol framing, versioning, backpressure, and error
  semantics.
\item
  Validate throughput/latency behavior in distributed and multi-client
  scenarios.
\end{itemize}

\subsubsection{Phase 4: NUMA-focused optimization and
publication}\label{phase-4-numa-focused-optimization-and-publication}

\begin{itemize}
\tightlist
\item
  Implement and compare NUMA-aware allocation/scheduling strategies.
\item
  Publish controlled benchmark comparisons across topology and policy
  variants.
\item
  Finalize white paper v1.0 with reproducible evidence.
\end{itemize}

\subsection{12. Conclusion}\label{conclusion}

Furrballs aims to become a practical, locality-aware platform for
high-performance workloads where predictable latency matters more than
broad feature breadth. The project now follows an explicit path: ship a
strong Core first, expand library capabilities second, then build
Server/Client protocol layers that wrap and expose Core safely. This
phased strategy preserves delivery velocity while keeping the long-term
NUMA-aware objective central.

\subsection{Appendix A: Results Table
Template}\label{appendix-a-results-table-template}

{\def\LTcaptype{none} % do not increment counter
\begin{longtable}[]{@{}
  >{\raggedright\arraybackslash}p{(\linewidth - 16\tabcolsep) * \real{0.0938}}
  >{\raggedleft\arraybackslash}p{(\linewidth - 16\tabcolsep) * \real{0.1250}}
  >{\raggedright\arraybackslash}p{(\linewidth - 16\tabcolsep) * \real{0.0938}}
  >{\raggedright\arraybackslash}p{(\linewidth - 16\tabcolsep) * \real{0.0938}}
  >{\raggedleft\arraybackslash}p{(\linewidth - 16\tabcolsep) * \real{0.1250}}
  >{\raggedleft\arraybackslash}p{(\linewidth - 16\tabcolsep) * \real{0.1250}}
  >{\raggedleft\arraybackslash}p{(\linewidth - 16\tabcolsep) * \real{0.1250}}
  >{\raggedleft\arraybackslash}p{(\linewidth - 16\tabcolsep) * \real{0.1250}}
  >{\raggedright\arraybackslash}p{(\linewidth - 16\tabcolsep) * \real{0.0938}}@{}}
\toprule\noalign{}
\begin{minipage}[b]{\linewidth}\raggedright
Scenario
\end{minipage} & \begin{minipage}[b]{\linewidth}\raggedleft
Threads
\end{minipage} & \begin{minipage}[b]{\linewidth}\raggedright
Compression
\end{minipage} & \begin{minipage}[b]{\linewidth}\raggedright
NUMA Policy
\end{minipage} & \begin{minipage}[b]{\linewidth}\raggedleft
Read p99 (ms)
\end{minipage} & \begin{minipage}[b]{\linewidth}\raggedleft
Write p99 (ms)
\end{minipage} & \begin{minipage}[b]{\linewidth}\raggedleft
Throughput (ops/s)
\end{minipage} & \begin{minipage}[b]{\linewidth}\raggedleft
Hit Ratio
\end{minipage} & \begin{minipage}[b]{\linewidth}\raggedright
Notes
\end{minipage} \\
\midrule\noalign{}
\endhead
\bottomrule\noalign{}
\endlastfoot
Baseline & & & & & & & & \\
Variant A & & & & & & & & \\
Variant B & & & & & & & & \\
\end{longtable}
}

\subsection{Appendix B: Environment
Template}\label{appendix-b-environment-template}

\begin{itemize}
\tightlist
\item
  CPU model and socket count:
\item
  NUMA node topology:
\item
  RAM size and speed:
\item
  Kernel and OS version:
\item
  Compiler and flags:
\item
  Build type:
\item
  Dataset profile:
\item
  Commit identifier:
\end{itemize}




\end{document}
